%% ma: L12-B14-0140
%% lop: 12
%% chuong: 5
%% bai: 14
%% dang: 12.14.5
%% loai: DS
%% muc: [NB, TH, NB, TH]
%% dap_an: "ĐĐĐS"
%% nguon: "Ảnh giáo viên gửi 25/11/2026 (tài liệu ôn tập Phương trình mặt phẳng), Phần II Câu 3 (điểm A(1;2;3), B(-1;0;5))"
%% kien_thuc: "Bài 7 (tọa độ vectơ, trung điểm), Bài 14 (phương trình mặt phẳng)"
%% trang_thai: chinh-thuc
%% ghi_chu: "Đáp án tự giải (sympy). Ý d) sai vì mặt phẳng trung trực là x+y-z+3=0."
\begin{ex}
Trong không gian $Oxyz$, cho hai điểm $A(1;2;3)$ và $B(-1;0;5)$.
\choiceTF
{\True $\overrightarrow{AB}=(-2;-2;2)$.}
{\True Mặt phẳng qua $A$ và vuông góc với $AB$ có phương trình là $x+y-z=0$.}
{\True Trung điểm của đoạn thẳng $AB$ là điểm $I(0;1;4)$.}
{Mặt phẳng trung trực của đoạn thẳng $AB$ có phương trình là $x+y-z+4=0$.}
\loigiai{a) Đúng.
b) $\overrightarrow{AB}=-2(1;1;-1)$, mặt phẳng qua $A$ nhận $(1;1;-1)$ làm vectơ pháp tuyến: $(x-1)+(y-2)-(z-3)=0\Leftrightarrow x+y-z=0$. Đúng.
c) $I\left(\dfrac{1-1}2;\dfrac{2+0}2;\dfrac{3+5}2\right)=(0;1;4)$. Đúng.
d) Mặt phẳng trung trực đi qua $I(0;1;4)$, nhận $(1;1;-1)$ làm vectơ pháp tuyến: $x+(y-1)-(z-4)=0\Leftrightarrow x+y-z+3=0$. Sai.}
\end{ex}
